% Supplement: explicit semantics, derivations, countermodels, and retrieval practice.
\chapter{严格推导补篇：证明、反模型与解释分歧}

\section{把“推出”写成可以检查的定义}
前面二十五讲提供法律主题的模型入口。本章补上“严格”二字需要的语义基础。以下定义用于核查推理，不能取代法条、证据规则或者解释论证。

固定一个具体问题，设有限原子命题集合为$\mathcal P=\{p_1,\ldots,p_n\}$。原子命题应足够具体，例如“甲在实施行为时已满16周岁”，不能只写“甲符合条件”。一个\textbf{赋值}$v$将每个原子映射为$0$或$1$，再按通常命题逻辑解释合取、析取、否定与蕴涵。

把法律前提$N$、已明确选择的解释方案$I$、已经给定的事实$F$合在一起，记为$\Gamma=N\cup I\cup F$。定义
\[
 \mathcal M(\Gamma)=\{v:\text{$v$使$\Gamma$中每个命题均为真}\}.
\]
这是该组前提仍容许的模型集合。由此定义
\[
 \Gamma\models Q
 \quad\Longleftrightarrow\quad
 \forall v\in\mathcal M(\Gamma),\ v(Q)=1.
\]
意思是：在每一种符合这些前提的赋值中，结论$Q$都成立。实际写论证时还应要求$\mathcal M(\Gamma)\ne\varnothing$，并注明每项事实的证据来源。若前提自相矛盾，不能利用空集合上的全称命题给任意定罪结论盖章。

\begin{modelprop}[反模型足以否定逻辑蕴涵]
若存在一个$v\in\mathcal M(\Gamma)$，使$v(Q)=0$，则$\Gamma\not\models Q$。
\end{modelprop}
\begin{proof}
由定义，$\Gamma\models Q$要求集合中每一个赋值均使$Q$为真。所给$v$属于该集合，却使$Q$为假，因此全称条件不成立。
\end{proof}

\caseinput{前提只有“若成立故意杀人既遂，则存在死亡结果”，另知“存在死亡结果”。拟推出“成立故意杀人既遂”。}
令$K$为故意杀人既遂，$D$为死亡结果。前提为$\{K\Rightarrow D,D\}$。取$v(K)=0,v(D)=1$，两项前提均为真，结论为假。故结论不由这些前提推出。过失致死、自然死亡等可以帮助理解缺口，但反模型本身已经足以证明此处的形式错误；具体哪种法律评价成立，仍须另行审查。

\memorycard{遇到“只要有结果，就一定有某罪”，将待检验主张写成$D\Rightarrow K$，与已知规则$K\Rightarrow D$分开。尝试$K=0,D=1$，说明为何它满足已知前提却否定拟推出的结论。}

\section{不完整信息：不能把问号偷偷改成零}
固定同一组规范$N$和解释$I$，设$F$只给定部分事实，记$\Gamma_F=N\cup I\cup F$。$\mathcal M(\Gamma_F)$中通常有多个赋值，可以把结论分类为
\[
 \operatorname{Status}_F(Q)=
 \begin{cases}
  1,&\forall v\in\mathcal M(\Gamma_F),\ v(Q)=1,\\
  0,&\forall v\in\mathcal M(\Gamma_F),\ v(Q)=0,\\
  ?,&\text{上述两种情况均不成立。}
 \end{cases}
\]
这里先要求$\mathcal M(\Gamma_F)\ne\varnothing$。这种“考察所有一致补全”的方法，能够保留同一个原子在不同位置出现时的相关性。

\begin{modelprop}[补充一致事实只会缩小模型集合]
若$F\subseteq F'$，则$\mathcal M(\Gamma_{F'})\subseteq\mathcal M(\Gamma_F)$。若两者均一致且$\Gamma_F\models Q$，则$\Gamma_{F'}\models Q$。
\end{modelprop}
\begin{proof}
任取$v\in\mathcal M(\Gamma_{F'})$。由于规范与解释保持不变，且$F$中每项也在$F'$中，$v$必满足$\Gamma_F$，故$v\in\mathcal M(\Gamma_F)$。若$Q$在较大的$\mathcal M(\Gamma_F)$中处处为真，在其子集上也处处为真。
\end{proof}

适用这条命题须固定同一组规则和解释。发现新证据后撤销了一个先前认定，属于\emph{修正前提}，不一定满足$F\subseteq F'$；新法改变规则，也不是单纯增加事实。因而不能据此宣称“案件知道得越多，原有有罪结论就越不会改变”。

\subsection{三值表的准确边界}
本册开头的三值合取表适合检查必要条件缺口，但不能不加限定地替代所有公式的经典语义。例如，尚不知$P$真假时，逐格代入强三值运算会把$P\vee\neg P$记为“?”；按一致补全观察，每个补全都使该式为真。两个方法回答的细节不同。含重复原子、互斥条件或关联约束时，宜直接考察完整模型，避免把相关的未知项当作可以任意独立变化。

\caseinput{甲是否自动放弃尚不明确，其他中止条件已给定。}
一组一致补全使自动性成立，另一组使它不成立，因此不能证明中止结论。这里既没有证明甲自动，也没有证明甲不自动。应回到停止原因、行为人的认识和外部条件，而不是用一个问号结束事实调查。

\section{必要、充分、等价：逐一提供证明义务}
设$C$为一个明确限定的法律结论，$P$为某项条件。
\begin{center}
\begin{tabular}{p{.18\textwidth}p{.28\textwidth}p{.43\textwidth}}
\toprule
声称 & 必须证明 & 一个否定它的反模型\\\midrule
$P$是必要条件 & $C\Rightarrow P$ & $C=1,P=0$\\
$P$是充分条件 & $P\Rightarrow C$ & $P=1,C=0$\\
$P$与$C$等价 & 两个方向分别成立 & 上述任一反模型\\
$P$只提供线索 & 不承诺充分或必要 & 仍须列出待补条件\\\bottomrule
\end{tabular}
\end{center}

若把一个规则写作$C\leftrightarrow P_1\wedge\cdots\wedge P_n$，就同时声称“列出的每一项均必要”以及“这些项合在一起已充分”。这比只写$C\Rightarrow P_j$强得多。许多刑法口诀只适合做检索线索，不能自动升级为双向等价式。

\begin{modelprop}[否定一个析取分支不等于否定整体]
设背景前提$\Gamma$一致，$Q=A\vee B_1\vee\cdots\vee B_m$，且$m\ge1$。若对某个$j$，$\Gamma\cup\{\neg A,B_j\}$仍一致，则$\Gamma\cup\{\neg A\}\not\models\neg Q$；若增加任一$B_j$为真这一前提，则可推出$Q$。
\end{modelprop}
\begin{proof}
由相容性前提，可取一个同时满足$\Gamma$、$\neg A$与$B_j$的赋值。它使$Q$为真，因而使$\neg Q$为假，是第一项推断的反模型。第二项依析取引入：任何一个分支为真，整个析取为真。这里不假定各分支任意独立；若背景规则已排除所有替代分支，就不能再使用这个反模型。
\end{proof}

盗窃罪的数额路径与多次、入户、携带凶器、扒窃等路径正适合用此命题检查。\lawref{L1}{第264条}\lawref{L7}{第1至3条} 证明的是路径表达式的逻辑性质；各分支是否实际满足及是否具备其他构成条件，不能省略。

\section{最小充分条件集：把背诵变成可拆解结构}
给定一致的有限条件集$\Gamma$与目标$Q$，若$S\subseteq\Gamma$满足$S\models Q$，称其为一个支持集。一致性保证这里没有利用矛盾前提的空真来制造支持。若对每个$s\in S$都有$S\setminus\{s\}\not\models Q$，则称$S$为\textbf{包含意义上的最小支持集}。这里的“最小”不表示所有支持集中元素最少，也不自动表示法律上不存在其他规则路径。

\begin{modelprop}[纯合取的逐项缺失反例]
设$p_1,\ldots,p_n$彼此独立，$C=p_1\wedge\cdots\wedge p_n$。事实集$S=\{p_1,\ldots,p_n\}$是$C$的最小支持集。
\end{modelprop}
\begin{proof}
首先，$S$使每个$p_i$为真，依合取引入得到$C$。其次，任删$p_j$，都可令其为假并令其余$p_i$为真；这个赋值满足剩余事实，却不满足$C$。故删除任一项后均不能推出$C$。
\end{proof}

独立性假设不能忽略：若规范或事实已蕴含$p_1\Rightarrow p_2$，则$p_2$可能不是这份证明记录中必须单独给出的事实，虽然它仍是法律构成的一项必要条件。学习时应区分\emph{法律要件}与\emph{证明该要件所需的事实材料}。

\memorycard{一张规则卡背面写全条件后，为每个条件各做一个“仅删此项”的变式。若删掉某项仍能沿另一规则路径得到同一结论，就记录那条路径，不强行宣称条件绝对必要。}

\section{争议学说作为参数：分别推导再比较}
解释争议不能通过偷偷选一个定义而消失。设两种方案分别为$I_1,I_2$，得到
\[
 \mathcal C_1=\{Q:N\cup I_1\cup F\models Q\},\qquad
 \mathcal C_2=\{Q:N\cup I_2\cup F\models Q\}.
\]
交集$\mathcal C_1\cap\mathcal C_2$表示两方案共同支持的结论；对称差
\[
 \mathcal C_1\mathbin{\triangle}\mathcal C_2
 =(\mathcal C_1\setminus\mathcal C_2)\cup(\mathcal C_2\setminus\mathcal C_1)
\]
定位方案导致的结论差异。此运算不能告诉我们哪种解释更有法律依据；选择方案仍须比较法条、体系、规范目的及相应学说论证。

\caseinput{一个取财案件中，其他条件固定，争点只在交付者交付财物时是否具有相应处分认识。}
先把两方案对处分认识的要求各写成一个谓词。再分别代入同一事实，不改变其他变量。若一个方案认可诈骗链中的处分环节，另一方案不认可，差异应定位到这一环节；不能把整道题说成“任何答案都可以”，也不能为使结论一致而临时修改事实。

\begin{modelprop}[共同结论对所列方案的稳定性]
若所列每一组前提$N\cup I_i\cup F$均一致并均推出$Q$，则$Q$在这组方案范围内稳定。仅凭这一有限比较，不能一般性地保证所有未列方案也支持$Q$。
\end{modelprop}
\begin{proof}
第一部分是有限全称命题的直接展开。为否定第二部分所述的一般性保证，取一个原子$Q$，令$N=F=\varnothing$，所列每个$I_i=\{Q\}$，另设$I_*=\{\neg Q\}$。各组前提均一致，已列方案都推出$Q$，而$I_*$不推出$Q$。这是一般推断的反例，并非声称对任意具体$Q$都能找到相反方案；若$N\cup F$本身已经推出$Q$，那么所有一致的补充解释都会保留该结论。
\end{proof}

\section{从推导树到可执行检查}
一份完整记录至少保留三类节点：\textbf{依据节点}写明条文或解释方案，\textbf{事实节点}写明证据或题目假定，\textbf{推论节点}写明所用规则与上游节点。下面画的是论证依赖关系，不是原书图的替代品。

\begingroup\setlength{\topsep}{3pt}
\begin{center}
\begin{tikzpicture}[>=Stealth,every node/.style={draw,align=center,font=\small,inner sep=5pt}]
\node (r) at (0,0) {规范前提\\$A\wedge B\Rightarrow C$};
\node (a) at (-3.5,-1.3) {事实节点\\$A$};
\node (b) at (3.5,-1.3) {事实节点\\$B$};
\node (ab) at (0,-1.3) {合取引入\\$A\wedge B$};
\node (c) at (0,-2.6) {肯定前件\\$C$};
\draw[->](a)--(ab);\draw[->](b)--(ab);
\draw[->](ab)--(c);\draw[->](r.east)--(5.1,0)--(5.1,-2.6)--(c.east);
\end{tikzpicture}
\end{center}
\endgroup

检查程序可以枚举小型布尔模型、核算金额边界、寻找反模型，或者核对每个推论是否有父节点。程序不能仅凭上述图确定$A$是否有证据支持，也不能证明引用的规则是当前适用的完整法律规范。本工程提供的模型检查脚本属于这种有限检查，不是全部刑法已经机器证明的声明。

\chapter{完整演算补篇：边界、归责与复习路径}

\section{正当防卫：先确定分支，再做真值运算}
本节对应第九讲。先固定行为具有防卫性质、不属于滥用防卫权等排除情形，再判断是否符合特殊防卫的全部法定条件。一般防卫过当要求“明显超过必要限度”和“造成重大损害”同时成立。\lawref{L1}{第20条}\lawref{L4}{限度与特殊防卫}

令$B$表示已通过前述基础审查，$S$表示符合特殊防卫条件；在一般防卫分支中令$X$为明显超过必要限度，$H$为造成重大损害。仅为组织检索，定义四个分支：
\begin{align*}
 D_0&=\neg B,\\
 D_1&=B\wedge S,\\
 D_2&=B\wedge\neg S\wedge X\wedge H,\\
 D_3&=B\wedge\neg S\wedge\neg(X\wedge H).
\end{align*}
$D_0$要求继续按其他规则评价；$D_1$进入特殊防卫审查后的结论；$D_2$满足所建模的一般防卫过当双条件；$D_3$表示这两个条件未同时具备。特别地，$D_3$不能被随意命名为“行为绝未超过限度”，因为它包括$X=1,H=0$的赋值。

\begin{modelprop}[防卫检索分支的互斥与穷尽]
在$B,S,X,H$均已取确定真假值的模型内，$D_0,D_1,D_2,D_3$恰有一个成立。
\end{modelprop}
\begin{proof}
若$B=0$，只有$D_0$成立。若$B=1$而$S=1$，只有$D_1$成立。若$B=1,S=0$，再按$X\wedge H$的真假分类，分别只有$D_2$或$D_3$成立。各对分支在$B$、$S$或$X\wedge H$至少一项上要求相反，故不会同时成立。
\end{proof}

\caseinput{题目已确认基础条件$B=1$、不属于特殊防卫$S=0$；重大损害$H=1$，但$X$没有证据结论。}
第一步，排除$D_0,D_1$。第二步，保留两个一致补全：$X=1$落入$D_2$，$X=0$落入$D_3$。第三步，因结论分支不同，不能据$H=1$完成防卫过当的推导。第四步，补查侵害的性质、强度、紧迫性、双方力量等与限度有关的事实，不能直接用损害严重程度替代$X$。

\memorycard{先背分支前提$B,S$，再背双条件$X,H$。闭卷画出四分支，并回答：为什么$X=1,H=0$不能被语言误写成“没有超过限度”？}

\section{犯罪形态：时间、终点和停止原因不能压成一个轴}
本节对应第十一讲。给定罪名$k$，将其特定既遂标准写为$F_k$，着手标准写为$S_k$，不得用一个统一的“财产到账”或“动作做完”替代所有罪名。\lawref{L1}{第22至24条}

本式只用于故意犯罪既遂前的犯罪过程中，行为人已放弃本次犯罪（非暂时停顿）的情形；既遂后的返还、救助不属本式适用域。令$V$为自动性，$R$为当时还需要有效防止结果，$E$为所要求的有效防止。在这个已限定的域内，自动中止路径的检索条件可写成
\[
 Z=V\wedge(\neg R\vee E).
\]
这一式子把“有的情形仅停止实行即可”“有的情形必须有效防止结果”分开。$R$的判断依犯罪进程与具体规范确定，不能由行为人自行宣布。

\begin{modelprop}[需要防止结果时，仅有后悔不充分]
在$R=1$的分支中，$Z$化简为$V\wedge E$。因此$V$本身不是$Z$的充分条件。
\end{modelprop}
\begin{proof}
代入$R=1$，得$\neg R\vee E=0\vee E=E$，故$Z=V\wedge E$。赋值$V=1,E=0$使$V$成立而$Z$不成立，即为充分性命题的反模型。
\end{proof}

在本式适用域内，也不能从$Z=0$直接推出未遂；还须确认已着手、未得逞、意志以外原因等相应条件。域尚未确定时，应先区分信息不足、行为尚未停止、已经达到法定既遂等情形，不能一概给它们赋值$Z=0$。对全体案件把“既遂、中止、未遂、预备”写成未经限定的四选一函数，会掩盖这些前提。

记忆时画两个互相配合的表：横轴表记录准备、着手、实行终了、既遂等实际时点，不假定所有罪名均按同一顺序出现；条件表记录自动性、停止原因、是否需要防止结果及有效性。每次只改变一个变量，观察评价变化，避免把时间更晚误当成处罚当然更重。

\section{共同犯罪：逐人逐项构造归责矩阵}
本节对应第十二讲。设参与者集合$P=\{p_1,\ldots,p_m\}$，待归责行为或结果集合$R=\{r_1,\ldots,r_n\}$。分别建立
\[
 A_{ij}=\text{存在可论证的行为贡献},\quad
 I_{ij}=\text{属于该人的相应共同故意范围},\quad
 C_{ij}=\text{符合其他归责与责任条件}.
\]
这里仅检查须由共同故意覆盖的具体行为或结果；若某项加重结果依法允许过失归责，就必须另设相应主观条件，不能原样套用这一矩阵。在此范围内，故意共犯归责候选为$G_{ij}=A_{ij}\wedge I_{ij}\wedge C_{ij}$。此矩阵把个人条件分开，不能把一行复制给所有人。共同过失不以共同犯罪论处等规则仍须独立适用。\lawref{L1}{第25至29条}

\begin{modelprop}[共同范围不同，归责行可以不同]
即使两个参与者对同一结果均有行为贡献，只要其中一人的相应故意必要条件为$0$，两人的故意共犯候选结论就不能由行为贡献相同而被强行视为相同。
\end{modelprop}
\begin{proof}
给出反模型：设$A_{1j}=A_{2j}=1$，$I_{1j}=1,I_{2j}=0$，并取$C_{1j}=1$。则$G_{1j}=1$；无论$C_{2j}$为何，$G_{2j}=1\wedge0\wedge C_{2j}=0$。故行为贡献列相同不保证归责列相同。这不表示两行必然不同：若$C_{1j}=0$，两项候选也可能同为$0$。
\end{proof}

\caseinput{甲乙共谋盗窃，乙望风；甲另起杀人故意实施杀人，乙对此不知情，也无针对杀人的帮助。}
建两列“盗窃”“杀人”，建两行“甲”“乙”。先填共同盗窃的参与和故意，再单列甲的超出行为。对乙的杀人列不能从“共谋盗窃”补出$I_{2,\text{杀人}}=1$。若题目另增加乙事前支持杀人或当场有意加入等事实，应重新审查该列，而非沿用旧输入。

\modelbound{矩阵没有解决片面共犯、承继共犯及身份犯等全部解释争议；这些问题仍需固定$I$后分别计算。矩阵的价值是迫使论证写明“谁对什么负责”，而不是自动确定每个人罪名。}

\section{数罪并罚：上下界可算，具体刑期不能只靠加法}
本节对应第十五讲。仅讨论已确定的若干有期徒刑在刑法第69条相应规则内的并罚区间；其他刑种、特别规定及法定情节另行处理。\lawref{L1}{第69至71条}
设各罪刑期为$t_1,\ldots,t_n>0$，令$S=\sum_i t_i$、$M=\max_i t_i$，上限函数为
\[
 U(S)=\begin{cases}20,&S<35,\\25,&S\ge35.\end{cases}
 \qquad
 \mathcal Y=\bigl[M,\min\{S,U(S)\}\bigr].
\]
$\mathcal Y$记录本模型给出的刑期数值上下界，不表示其中每一个实数都可以作为合法宣告刑；具体计算单位及其他适用要求仍须核对。实际执行刑期$y$仍需依法裁量，不能从该区间凭数学唯一确定一个数。

\begin{modelprop}[区间非空的准确条件]
对给定数值，$\mathcal Y\ne\varnothing$当且仅当$M\le\min\{S,U(S)\}$。由于$M\le S$，该条件等价于$M\le U(S)$。
\end{modelprop}
\begin{proof}
实数闭区间$[a,b]$非空当且仅当$a\le b$，代入即可。$M\le S$由所有$t_i>0$得到，因此还需检查的限制只有$M\le U(S)$。
\end{proof}

若一个输入让下界大于上界，应先检查是否把不属于这一简单模型的刑种、既有并罚刑或特别规则硬塞进来，不得把空区间“修正”为任意期望刑期。模型适用性优先于代数运算。

\caseinput{三个已判刑期分别为12、10、8年。}
$S=30,M=12,U(S)=20$，因此$12\le y\le20$。若改成15、12、8年，$S=35,M=15,U(S)=25$，因此$15\le y\le25$。跨越35年这个分界改变的是上限函数，不能把所有情形统一套成20年。

\subsection{先并后减与先减后并为什么不能交换}
设$B(a,b)$只表示在具体规范与裁量下确定的并罚结果。漏罪与新罪适用的法律次序不同；例如符号形式$B(a,b)-e$与$B(a-e,b)$中，$e$是已经执行部分。即使暂取一个简化数值算子$B(a,b)=\max(a,b)$用于展示顺序，$a=10,b=8,e=3$时两式分别为$7$和$8$。这只是“不应默认两种运算可交换”的反例，\textbf{不是实际并罚算法}；真正计算必须依第70、71条及案件条件。

\memorycard{刑期卡先写适用范围，再写$S,M,U(S)$，最后才算区间。次序卡各画一条箭头，并在箭头上写对应法条；不能只背“加减顺序”而不认案件是漏罪还是新罪。}

\section{财产犯罪：状态转换与取得机制分别建模}
本节对应第十九讲。设财物$o$在时点$t$的占有关系为$\operatorname{Pos}(p,o,t)$。占有是具有规范评价内容的关系，不等同于GPS坐标；所有权记为$\operatorname{Own}$，处分地位记为$\operatorname{Disp}$，三者不得混用。

一次取财分析写成
\[
 (s_0,a,m,s_1,u),
\]
其中$s_0,s_1$分别记录行为前后的权利与占有状态，$a$是行为，$m$是使状态改变的机制，$u$是相应主观条件。仅保留起点与终点会丢失$m$；盗窃、诈骗、抢劫等都可能出现财物从乙转入甲控制的外观，区别必须回到取得机制及完整构成。\lawref{L1}{第263、264、266、270、271条}

\begin{modelprop}[相同终点不能唯一识别取得机制]
设两个机制$m_1\ne m_2$均可把$s_0$变为$s_1$。仅由有序对$(s_0,s_1)$不能在两者之间唯一确定机制。
\end{modelprop}
\begin{proof}
若存在仅依赖$(s_0,s_1)$且总能正确识别机制的函数$f$，它面对两种过程都必须给出同一输出；正确性却分别要求输出$m_1$与$m_2$。因$m_1\ne m_2$而矛盾。
\end{proof}

\caseinput{过程甲：行为人平和地秘密拿走他人控制中的手机；过程乙：行为人欺骗占有人，使其陷入错误并处分同一手机。其余相关条件按各自构成另行补齐。}
两过程的最后控制状态可相同，但取得机制不同。因此“手机最后在谁手里”只能回答结果问题，不能完成盗窃与诈骗的区分。

\subsection{三角诈骗的关系表}
把欺骗者$d$、受骗人$p$、财产受损者$b$分开；本小节用$b$表示人，避免与前述财物$o$混淆。关系表至少有：欺骗$d\to p$、错误$p$、处分$p\to\text{财产}$、处分地位$p\rightsquigarrow b$、损失$b$。在采用要求相应处分地位的解释方案时，一项机制检查写为
\begin{align*}
 K_I(d,p,b)={}&D(d,p)\wedge E(p)\wedge P(p,b)\wedge A_I(p,b)\wedge L(b)\\
 &\wedge C_{DE}(d,p)\wedge C_{EP}(p,b)\wedge C_{PL}(p,b).
\end{align*}
其中$P(p,b)$表示$p$处分归属于受损者$b$的财产；$C_{DE}$、$C_{EP}$、$C_{PL}$分别要求有关欺骗引起有关错误、该错误引起该处分、该处分造成相应财产损失。仅把这些事件同时列为真，并不能替代因果与规范联系的论证。本式只检索这条机制，并非诈骗罪完整构成，仍须补足相应非法占有目的等条件。

$A_I$需按所选解释论证。没有它，不能因受骗人“帮忙把东西递出来”就当然认定三角诈骗既遂；也不能不经论证就把所有第三人交付视为盗窃。

\begin{modelprop}[关系连接必须保留同一主体]
从$D(d,p)$与$P(q,b)$不能直接连接出“$p$受骗后处分$b$财产”的链，除非另有$p=q$或足以连接两人的具体机制。
\end{modelprop}
\begin{proof}
取$p\ne q$，令$p$确实受骗而未处分，令$q$基于无关原因处分。两项原前提都成立，但拟推出的连续机制不存在。因此缺少连接前提时推论不成立。
\end{proof}

即使补入$p=q$，也只解决主体同一问题；仍要证明错误确由有关欺骗引起、处分确由该错误引起。对多人收款、公司付款、自动设备取财等题目，每根箭头写明主体和机制，能够发现把甲的错误认识、乙的交付和丙的损失随意拼接的问题。

\subsection{金额的类型约束}
将犯罪数额、损失、获利、追缴、赔偿分别记为带标签的量
\[
 v_c:\mathsf{CrimeAmount},\quad v_l:\mathsf{Loss},\quad
 v_g:\mathsf{Gain},\quad v_r:\mathsf{Recovery},\quad v_d:\mathsf{Compensation}.
\]
它们可能数值相等，但相等必须有依据。例如退赔使某项未弥补损失减少，不当然使原行为阶段、犯罪数额或既遂事实倒退消失。若某条规范规定退赔影响某种处罚后果，应把该规则单独接入，不能私设恒等式$v_c=v_l-v_d$适用于全部罪名。

\memorycard{一张财产题卡只准使用五栏：“原占有—取得机制—处分主体与地位—新控制—主观目的”。金额另外开五个格，不许用一个“钱数”包办所有判断。}

\section{职务犯罪：身份标签必须连接到具体职务关系}
本节对应第二十三、二十四讲。把“某人有某身份”与“其利用的正是该职务便利”分开。设$O(p,j,t)$表示人在时点$t$任职$j$，$A(j,a)$表示该职务对事项$a$具有相应权限、管理或职责联系，$U(p,j,a)$表示行为人实际利用这条联系。仅有$O$不能推出$U$。

\begin{modelprop}[任职事实不能独立证明职务利用]
从$O(p,j,t)$本身不能推出$U(p,j,a)$。
\end{modelprop}
\begin{proof}
取一个人在该时点任职，但案件事项与其职务无所要求的联系，或者其行为并未利用职务便利的模型。$O$为真而$U$为假，构成反模型。
\end{proof}

这是区分职务便利、工作接近机会、私人关系与影响力路径的基本检查。具体何种联系满足贪污、受贿、斡旋受贿或利用影响力受贿等构成，应分别根据法条展开，不能把“认识官员”当作统一充分条件。\lawref{L1}{第382、385、388、388条之一}

\subsection{中间联系的不同路径不能合并}
以请托人$q$、中间人$m$、实际办理事项者$s$为三类角色。先写$q\to m$的利益给付或请托关系，再写$m\to s$的联系性质，最后写$s\to a$的职务行为。$m\to s$可能属于职权或地位形成的便利，也可能属于近亲属、关系密切者或离职影响力等法定路径；必须保留构成所需身份、时间、便利、利益性质与共同故意条件。只有结构看似同为三个人，不足以把不同罪名视为同一个模型。

\subsection{数额阈值中的析取演算}
2026年解释（二）第一条对单位受贿“情节严重”所列的一组金额路径，在以万元为单位并限定相应适用范围后，可写成
\[
 Q(v,E)=(v\ge20)\vee(10\le v<20\wedge E),
\]
$E$表示具备该条列明的相关情形。\lawref{L8}{第1条} 这是该项门槛的抽取，不是单位受贿罪完整构成。

当$v=15$时，$Q(15,E)=0\vee(1\wedge E)=E$；当$v=20$时，$Q(20,E)=1$，不再需要通过第二金额分支。把这一规则记成“统一20万元起”会遗漏$10\le v<20$且$E=1$的分支；记成“统一10万元起”又会把$E=0$的案件错误纳入。

\memorycard{职务题画三层：人的身份、与事项的职务联系、实际利用方式。金额题画区间，并在较低区间写额外情节变量；边界点10、20各算一次。}

\section{覆盖二十五讲的闭卷演算安排}
这张表把每讲的记忆任务指定为一个可检查的动作。动作完成只说明会重建核心结构，不能替代该讲全部法条和案例学习。

\begingroup\small\linespread{1}\selectfont
\setlength{\LTpre}{4pt}\setlength{\LTpost}{4pt}
\begin{longtable}{p{.09\textwidth}p{.40\textwidth}p{.39\textwidth}}
\toprule
讲次 & 闭卷必须写出的结构 & 用来检查的最小变式\\\midrule
\endfirsthead
\toprule 讲次 & 闭卷必须写出的结构 & 用来检查的最小变式\\\midrule\endhead
1 & 规范依据、解释、空间、时间四层 & 只改变新旧法的轻重关系\\
2 & 完整构成与阻却、责任的依赖图 & 只删去一个必要条件\\
3 & 法益、对象、结果三个不同变量 & 有对象损伤但无所称法益侵害\\
4 & 自然人、身份、单位归责入口 & 员工身份不变而行为归属改变\\
5 & 义务来源、能力、未履行、归责链 & 只改变现实履行能力\\
6 & 基本结果与加重结果分层 & 有较重结果但无所需归责\\
7 & 事实因果与规范归责两张表 & 只改变介入因素的作用机制\\
8 & 认识对象、希望或放任、注意义务 & 只改变对特定结果的认识态度\\
9 & 基础条件、特殊防卫、过当双条件 & 保持损害不变而改变限度条件\\
10 & 年龄区间和各区间附加条件 & 保持罪名不变而跨越年龄边界\\
11 & 阶段、自动性、必要有效性 & 只改变救助是否实际有效\\
12 & 每人一行、每结果一列的矩阵 & 加入一项超出共同故意的行为\\
13 & 行为单元、候选罪名、处理依据 & 同一结果数而行为单元改变\\
14 & 主刑、附加刑和非刑罚措施的类型 & 保留名称近似而更换法律性质\\
15 & 法定刑幅、情节、并罚运算次序 & 总和从34变35；漏罪改为新罪\\
16 & 执行状态和时效起算、停止问题 & 只改变法定最高刑或追诉事件\\
17 & 构成规定、注意规定、拟制入口 & 删除拟制条文后重算候选罪名\\
18 & 法益、行为方式、目的及身份 & 只改变控制他人的目的与机制\\
19 & 占有转换、取得机制、处分关系 & 终点不变而把欺骗改成秘密取走\\
20 & 特定对象风险与公共危险的区别 & 只改变危险可扩散的范围\\
21 & 行政前提与刑事构成两层 & 行政违法成立而严重性条件缺失\\
22 & 具体对象、行为、身份及法定程度 & 只改变诉讼角色或违禁品行为方式\\
23 & 身份、职务联系、利益、交易路径 & 只改变中间人与办理者的联系性质\\
24 & 职责、违背职责、结果、归责 & 有任职身份但无该项具体职责\\
25 & 具体行为、境外关系、对象与主观条件 & 只有境外联系而无所要求行为\\\bottomrule
\end{longtable}
\endgroup

每次检索按三轮进行：第一轮只写变量与规则，不翻答案；第二轮为漏掉的条件造一个反例；第三轮解释该条件为何改变推导。复习记录不只写对错，还应记下失败发生在哪一条边：没有记住规范、没有识别事实、把必要条件当充分条件、误连主体、忽略例外，或者套错版本。这样才能决定下一次应背法条、练案例还是重画推导图。

\modelbound{本补篇证明的对象是明确模型中的逻辑、集合与算术命题。规范是否成立、学说如何选择、事实是否证明，仍须分别论证。前提透明、推导可检查、反例可定位，是这里“严格数学化”的具体含义；不能把它扩张为法律判断已经成为纯计算。}
